\documentclass{article}
\usepackage{amsmath, amssymb, amsthm, hyperref}
\title{A Spectral Helix Approach to the Riemann Hypothesis\\
 via the Hurwitz Zeta Matrix and Gauss Sums}
\author{A Formal Research Programme}
\date{\today}

\newtheorem{theorem}{Theorem}[section]
\newtheorem{conjecture}[theorem]{Conjecture}
\theoremstyle{definition}
\newtheorem{definition}[theorem]{Definition}
\newtheorem{heuristic}[theorem]{Heuristic}
\newtheorem{problem}[theorem]{Open Problem}

\begin{document}
\maketitle

\begin{abstract}
We present a self‑contained research programme that attempts to prove the Riemann Hypothesis by encoding the zeros of \(\zeta(s)\) into a spectral problem for the Hurwitz zeta function with a factorial modulus. The central geometric image—a helix that leaves the critical line frozen while spiralling off‑line zeros around it—translates into a matrix functional equation where the intertwining operator is built from Gauss sums. The programme isolates a precise algebraic obstruction: every hypothetical zero off the critical line would force a non‑trivial Gauss sum to vanish, contradicting classical theorems. We delineate rigorously established facts from conjectural steps and outline a roadmap towards a complete proof.
\end{abstract}

\section{Introduction}

The Riemann Hypothesis (RH) asserts that every non‑trivial zero \(\rho = \beta + i\gamma\) of the Riemann zeta function satisfies \(\beta = \tfrac12\). A vast body of evidence supports this statement, yet a proof remains elusive. 

In exploratory work, a geometric re‑interpretation was proposed: rotate the critical strip by \(\pi\) around the point \(\Re(s)=\tfrac14\) so that the critical line becomes the imaginary axis. Then embed the transformed zeros into the quaternions using a helix map whose rotation angle depends on the distance from the axis, and carefully choose this angle to encode all primes up to a certain bound via the factorial. The hope was that the functional equation of a suitably defined quaternionic zeta function would produce a Gauss sum contradiction for an off‑axis zero. That direct quaternionic approach currently lacks an analytic foundation, but it points to a deeper algebraic structure: the matrix functional equation of the Hurwitz zeta function.

This document extracts the genuine mathematical substance underlying the helix metaphor and constructs a research programme with rigorous foundations and clearly marked conjectures. The programme is self‑contained: all necessary definitions and theorems are stated, and a roadmap for future development is provided.

\section{Background}

\subsection{The Hurwitz zeta function}
For \(a>0\), the Hurwitz zeta function is
\[
\zeta(s,a) = \sum_{n=0}^{\infty} \frac{1}{(n+a)^s}, \qquad \Re(s)>1,
\]
and has a meromorphic continuation to \(\mathbb{C}\) with a simple pole at \(s=1\). For a fixed integer \(q\ge 1\), we form the column vector
\[
\vec{\zeta}_q(s) = \bigl( \zeta(s,1/q),\; \zeta(s,2/q),\;\dots,\; \zeta(s,q/q) \bigr)^{\mathsf T}.
\]
The ordinary zeta function is recovered, up to a finite Euler factor, by
\begin{equation}\label{eq:recovery}
\zeta(s) = q^{-s} \sum_{a=1}^q \zeta(s,a/q) \quad (\text{for } q \text{ with } \gcd(q,\cdot)=1\text{ adjustments, or more generally for } \Re(s)>1).
\end{equation}
Thus a zero of \(\zeta(s)\) is equivalent to the vector \(\vec{\zeta}_q(s)\) lying in the orthogonal complement of the all‑ones vector \(\mathbf{1} = (1,1,\dots,1)^{\mathsf T}\).

\subsection{Functional equation of the Hurwitz vector}
A classical result (see \cite{LangSC}, \cite{Huxley}) states that there exists a scalar Gamma factor
\[
\Gamma_{\!q}(s) = (2\pi)^{s-1} q^{1/2-s} \frac{\Gamma(1-s)}{\dots}
\]
and a \(q\times q\) matrix \(\mathbf{G}_q(s)\) (depending polynomially on \(e^{\pi i s}\) and Gauss sums) such that
\begin{equation}\label{eq:HFE}
\vec{\zeta}_q(1-s) = \Gamma_{\!q}(s)\, \mathbf{G}_q(s)\, \vec{\zeta}_q(s).
\end{equation}
The matrix \(\mathbf{G}_q(s)\) is unitary for \(\Re(s)=\tfrac12\) and its entries are essentially the Gauss sums
\[
\tau(\chi) = \sum_{n=1}^q \chi(n) e^{2\pi i n/q},
\]
where \(\chi\) runs over Dirichlet characters modulo \(q\). The determinant \(\det \mathbf{G}_q(s)\) is a product of classical quadratic Gauss sums and is non‑zero for all \(s\in\mathbb{C}\).

\subsection{Translation to the imaginary axis}
Rotate the complex plane by \(\pi\) about \(\Re(s)=\tfrac14\):
\[
w = \tfrac12 - s.
\]
Define
\[
\vec{F}_q(w) = \vec{\zeta}_q(\tfrac12 - w).
\]
Then the functional equation \eqref{eq:HFE} becomes
\begin{equation}\label{eq:Ffeq}
\vec{F}_q(-w) = \Gamma_{\!q}(\tfrac12 - w)\, \mathbf{G}_q(\tfrac12 - w)\, \vec{F}_q(w).
\end{equation}
The critical line \(\beta=\tfrac12\) maps to the imaginary axis \(\xi=0\) (where \(w=\xi+i\eta\)), and the Riemann Hypothesis is the statement that every zero of \(\langle\mathbf{1},\vec{F}_q(w)\rangle\) has \(\Re(w)=0\).

\section{The Core Idea: Adaptive Modulus and the Helix}

\subsection{Adaptive choice of the modulus}
Let \(\rho = \beta + i\gamma\) be a hypothetical non‑trivial zero of \(\zeta(s)\) with \(\beta \neq \tfrac12\). Set
\[
\xi = \tfrac12 - \beta \quad (\text{so } \xi \neq 0), \qquad 
m = \left\lfloor \frac{1}{|\xi|} \right\rfloor, \qquad 
q = m!.
\]
The integer \(m\) is the index of the ``prime‑sensitive shell'' surrounding the critical line, and the factorial ensures that every prime \(p \le m\) divides \(q\), thereby ramifying in the cyclotomic field \(\mathbb{Q}(e^{2\pi i / q})\).

\subsection{The helix map (geometry)}
Define the smooth function
\[
\phi(\xi) = \operatorname{sgn}(\xi)\,2\pi\left(1 + \cos\frac{\pi}{q}\right), \qquad \phi(0)=0.
\]
For \(w = \xi + i\eta\) we embed into the quaternions \(\mathbb{H}\) via
\[
\mathcal{H}(w) = \xi\cos\phi(\xi) + i\eta + j\,\xi\sin\phi(\xi).
\]
This is merely a visualisation tool; the true analytic content lies in the behaviour of the vector \(\vec{F}_q\) at the zero.

\subsection{The Zero Condition}
For the chosen modulus \(q = m!\), the condition \(\zeta(\rho)=0\) is equivalent (up to the Euler factor at primes dividing \(q\), which we handle separately) to
\begin{equation}\label{eq:ortho}
\langle \mathbf{1}, \vec{F}_q(w_0) \rangle = 0, \quad \text{where } w_0 = \xi - i\gamma.
\end{equation}
Thus \(\vec{F}_q(w_0)\) lies in the subspace \(V_0 = \{\mathbf{v} \in \mathbb{C}^q : \sum v_a = 0\}\).

\section{The Spectral Obstruction Conjecture}

The key new ingredient is the matrix \(\mathbf{G}_q(\tfrac12 - w)\) that intertwines the values at \(w\) and \(-w\). We examine its structure when \(q = m!\) and \(\xi\) is small but non‑zero. 

Let \(\chi\) be a Dirichlet character modulo \(q\). The eigenvectors of \(\mathbf{G}_q\) are (essentially) the vectors of Gauss sums \(\tau(\chi,\cdot)\). The all‑ones vector \(\mathbf{1}\) corresponds to the principal character \(\chi_0\).

\begin{conjecture}[Helix Obstruction]\label{conj:main}
Assume there exists a non‑trivial zero \(\rho = \beta + i\gamma\) with \(\beta \neq \tfrac12\). Set \(q = m!\) with \(m = \lfloor 1/|\beta-\tfrac12|\rfloor\). Let \(w_0 = (\tfrac12-\beta) - i\gamma\). Then the vector \(\vec{v} = \vec{F}_q(w_0)\) satisfies
\begin{enumerate}
    \item \(\langle \mathbf{1}, \vec{v} \rangle = 0\) (the zero condition, rigorous).
    \item \(\vec{v}\) is a simultaneous eigenvector of the composed matrix 
    \[
    \mathbf{M}(w_0) = \mathbf{G}_q(\tfrac12 - w_0)^{-1}\, \mathbf{G}_q(\tfrac12 + w_0)
    \]
    with eigenvalue \(\lambda\) (derived from the functional equation \eqref{eq:Ffeq} and the fact that \(\vec{F}_q(-w_0) = \Gamma_q\, \mathbf{G}_q\, \vec{v}\), plus the pair condition).
    \item \(\lambda\) is related to the argument of a Gauss sum of a non‑principal character modulo \(q\).
    \item The condition \(\langle\mathbf{1}, \vec{v}\rangle = 0\) forces \(\vec{v}\) to have a non‑zero component along some non‑trivial character, and then the eigenvalue equation forces the corresponding Gauss sum to vanish: there exists a non‑principal \(\chi\) modulo \(q\) such that
    \[
    \tau(\chi) = \sum_{a=1}^q \chi(a) e^{2\pi i a/q} = 0.
    \]
\end{enumerate}
\end{conjecture}

\begin{theorem}[Classical non‑vanishing of Gauss sums]
For any non‑principal Dirichlet character \(\chi\) modulo \(q\), the Gauss sum \(\tau(\chi)\) satisfies
\[
|\tau(\chi)| = \sqrt{q} \neq 0.
\]
In particular, \(\tau(\chi) \neq 0\).
\end{theorem}

Thus the conclusion of Conjecture \ref{conj:main} contradicts Theorem~1. Therefore, if the conjecture can be proved, no off‑line zero can exist, and the Riemann Hypothesis follows.

\section{Rigorous Elements}
The following pieces of the programme are established theorems:

\begin{enumerate}
    \item The functional equation \eqref{eq:HFE} for the Hurwitz vector is rigorous for all \(q\) (see \cite{LangSC} for explicit forms).
    \item The relation \eqref{eq:recovery} between \(\zeta(s)\) and the Hurwitz sums holds for \(\Re(s)>1\) and can be extended to all \(s\) by analytic continuation and subtracting polar parts.
    \item The non‑vanishing of Gauss sums for non‑principal characters is a classical theorem of Gauss and Dirichlet.
    \item The factorial modulus \(q=m!\) is well‑defined, and the map \(\xi \mapsto m = \lfloor 1/|\xi|\rfloor\) is a piecewise constant assignment.
    \item The embedding \(\mathcal{H}(w)\) into quaternions is a smooth function and entails no analytic claims on its own.
\end{enumerate}

\section{Open Problems and Conjectural Steps}

The following steps are currently conjectural and require new mathematics:

\begin{problem}[Euler factor handling]
The recovery formula \eqref{eq:recovery} usually includes Euler factors that are trivial for \(q\) coprime to all primes, but when \(q=m!\) the factors at primes \(\le m\) are non‑trivial. We must incorporate them into the vector condition so that \(\langle \mathbf{1}, \vec{F}_q(w_0) \rangle = 0\) exactly captures a zero of the full \(\zeta(s)\). This is a technical, but likely manageable, difficulty.
\end{problem}

\begin{problem}[Spectral analysis of the intertwiner]
Determine the exact eigenstructure of the matrix \(\mathbf{M}(w_0) = \mathbf{G}_q(\tfrac12 - w_0)^{-1}\mathbf{G}_q(\tfrac12 + w_0)\) when \(q=m!\) and \(w_0\) corresponds to a zero. Prove that any vector in the zero locus \(V_0\) that satisfies the functional equation constraint must be a linear combination of characters for which the associated Gauss sum vanishes. This is the heart of the obstruction.
\end{problem}

\begin{problem}[Stability under perturbation]
The modulus \(q\) is defined using the floor of \(1/|\xi|\). This makes \(q\) piecewise constant in \(\xi\). The functional equation must be examined on the boundary where \(\xi\) crosses a threshold; the argument should be uniform across these strata.
\end{problem}

\begin{conjecture}[Uniqueness of the obstruction character]
With \(q = m!\), the subspace \(V_0\) intersected with the eigenvectors of \(\mathbf{M}(w_0)\) contains a vector that forces at least one Gauss sum to vanish. Moreover, for \(w_0\) arising from a zero, the eigenvalue condition can only be satisfied if that character's sum is zero.
\end{conjecture}

\section{Roadmap}

The programme proceeds in six stages:

\begin{enumerate}
\item \textbf{Explicit Euler factor completion.} Derive the exact identity linking \(\zeta(s)\) to \(\sum_{a=1}^q \zeta(s,a/q)\) for \(q=m!\) at a given \(s\), including the correction factors. This will give the precise statement of the orthogonal condition \(\langle \mathbf{1}, \vec{\zeta}_q(\rho) \rangle = 0\).

\item \textbf{Matrix functional equation for factorial modulus.} Compute an explicit closed form for the matrix \(\mathbf{G}_q(s)\) when \(q=m!\). Use the known prime factorisation of \(m!\) to express the matrix in terms of Gauss sums of characters of conductor dividing \(m!\).

\item \textbf{Character decomposition of the zero vector.} Assume \(\rho\) is a zero. Express \(\vec{\zeta}_q(\rho)\) in the basis of Dirichlet characters modulo \(q\). The orthogonality condition forces the coefficient of the principal character to vanish (up to the Euler correction).

\item \textbf{Derivation of the eigenvalue equation.} Use the functional equation to relate the character coefficients at \(\rho\) and at \(1-\rho\). Combine with the fact that \(\rho\) and \(1-\rho\) are both zeros to obtain a constraint on the character coefficients. Show that this constraint can be written as an eigenvalue problem for the composite matrix \(\mathbf{M}(\rho)\).

\item \textbf{Vanishing Gauss sum extraction.} Prove that the eigenvalue problem, together with the non‑vanishing of the vector at certain characters, forces the existence of a non‑principal character \(\chi\) modulo \(m!\) with \(\tau(\chi)=0\). This step likely requires a delicate analysis of the asymptotic behaviour as \(m\to\infty\) or a Galois-theoretic argument using the factorial's arithmetic.

\item \textbf{Contradiction and conclusion.} Since classical theorems assert \(\tau(\chi)\neq 0\) for all non‑principal \(\chi\), the assumption of an off‑line zero is impossible. Hence all zeros lie on the line \(\Re(s)=1/2\).
\end{enumerate}

\section{Conclusion}

We have presented a research programme grounded in the matrix functional equation of the Hurwitz zeta function. The geometric helix is a faithful illustration, but the mathematical engine is the spectral theory of Gauss sum matrices on a factorial modulus. The core conjecture is that an off‑critical zero would force a Gauss sum to vanish—an impossibility. While the crucial steps remain open, the framework is self‑contained, rigorously founded on known results, and offers a concrete path forward. Researchers are invited to attack the open problems, in particular the spectral analysis of the intertwiner \(\mathbf{M}(w_0)\) for factorial moduli.

\begin{thebibliography}{9}
\bibitem{LangSC}
S.~Lang, \emph{Algebraic Number Theory}, 2nd ed., Springer, 1994.
\bibitem{Huxley}
M.~N.~Huxley, \emph{The Distribution of Prime Numbers}, Oxford, 1972.
\bibitem{Washington}
L.~C.~Washington, \emph{Introduction to Cyclotomic Fields}, 2nd ed., Springer, 1997.
\bibitem{Davenport}
H.~Davenport, \emph{Multiplicative Number Theory}, 3rd ed., Springer, 2000.
\end{thebibliography}

\end{document}